\section{Os operadores}
\label{sec:operadores}

Esta seção descreve os quatro operadores tal como implementados. Cada um recebe a half-edge da aresta nova que parte do
vértice $v$ (chamada \code{edgeHE} ou \code{he} no código) e a half-edge do lado curto correspondente do novo elemento
(\code{elementHE}, $h_3$ ou $h_1$). Nos pseudocódigos, $\V[h]\leftarrow w$ corresponde a \code{setElementVertex(h, w)},
``$\Op[x]\leftrightarrow y$'' a \code{setOpposite(x, y)} com $y$ uma half-edge válida (liga nos dois sentidos) e
\code{duplicatePoint}$(v)$ cria um vértice novo, cópia de $v$, com as mesmas coordenadas.

As figuras desta seção foram geradas executando o código em malhas pequenas (\cref{app:repro}). Convenções visuais:
elementos contínuos em azul, elementos de interface em laranja, aresta a inserir em vermelho; vértices com aresta
colapsada têm contorno vermelho; vértices que pertencem a algum elemento de interface têm rótulo laranja; as cópias de
um vértice $X$ são chamadas $X_1, X_2,\dots$ na ordem de criação e são desenhadas ligeiramente afastadas umas das
outras. Arestas colapsadas (comprimento zero) aparecem como um pequeno segmento vermelho, com uma espessura
\emph{apenas visual}.

Todos os operadores usam a rotina de apoio \code{reindexElementsCCW} (\cref{alg:reindex}), que gira no sentido
anti-horário a partir de uma half-edge trocando o vértice de origem até encontrar uma half-edge de parada, a borda ou o
ponto de partida.

\begin{algorithm}[ht]
\caption{\code{reindexElementsCCW}$(h, \mathit{stop}, w)$}
\label{alg:reindex}
\begin{algorithmic}[1]
\State $s\gets h$
\Repeat
  \State $\V[h]\gets w$ \Comment{o canto de $h$ passa a usar o vértice $w$}
  \State $h\gets \Op[\prv(h)]$ \Comment{próximo elemento no sentido anti-horário}
\Until{$h=\mathit{stop}$ \textbf{ou} $h=\BORDER$ \textbf{ou} $h=s$}
\end{algorithmic}
\end{algorithm}

\subsection{Canônico}
\label{sec:canonico}

\textbf{Quando:} o vértice $v$ ainda não pertence a nenhum elemento de interface. É a ``inicialização'' do vértice.

\textbf{O que faz:} se $v$ é interno, nada é duplicado: o lado curto do novo elemento vira uma aresta colapsada de $v$
para $v$, com oposta \COLLAPSED, e ela é registrada em \code{\_collapsed[v]}. A trinca termina em $v$ e o elemento de
interface fica ``fechado'' ali. Se $v$ está na borda, a abertura chega ao exterior: cria-se um vértice $v'$, todos os
elementos do lado de \code{edgeHE} até a borda passam a usá-lo (rotação anti-horária), e o lado curto do novo elemento
fica na borda (\cref{alg:canonical,fig:canonico}).

\begin{algorithm}[ht]
\caption{\code{canonical}$(e, x)$ \quad ($e$ = \code{edgeHE}, $x$ = \code{elementHE})}
\label{alg:canonical}
\begin{algorithmic}[1]
\State $v\gets\V[e]$
\If{\code{isBorder}$(e)$} \Comment{gira em torno de $v$ procurando uma half-edge sem oposta}
  \State $v'\gets$ \code{duplicatePoint}$(v)$
  \State $\V[x]\gets v'$
  \State \code{reindexElementsCCW}$(e, \BORDER, v')$ \Comment{lado de $e$, até a borda}
  \State $\Op[x]\gets\BORDER$; \quad \code{\_inInterfaceElement}$[v']\gets$ \textbf{true}
\Else
  \State $\V[x]\gets v$; \quad $\Op[x]\gets\COLLAPSED$; \quad \code{\_collapsed}$[v]\gets x$
\EndIf
\State $\V[\nxt(x)]\gets v$; \quad \code{\_inInterfaceElement}$[v]\gets$ \textbf{true}
\end{algorithmic}
\end{algorithm}

\begin{figure}[ht]
\centering
\begin{subfigure}[b]{0.28\linewidth}\centering\begin{tikzpicture}[scale=1.400, every node/.style={font=\scriptsize}]
\filldraw[cefill] (0.000,0.000) -- (1.000,0.000) -- (1.000,1.000) -- (0.000,1.000) -- cycle;
\filldraw[cefill] (1.000,0.000) -- (2.000,0.000) -- (2.000,1.000) -- (1.000,1.000) -- cycle;
\filldraw[cefill] (2.000,0.000) -- (3.000,0.000) -- (3.000,1.000) -- (2.000,1.000) -- cycle;
\filldraw[cefill] (0.000,1.000) -- (1.000,1.000) -- (1.000,2.000) -- (0.000,2.000) -- cycle;
\filldraw[cefill] (1.000,1.000) -- (2.000,1.000) -- (2.000,2.000) -- (1.000,2.000) -- cycle;
\filldraw[cefill] (2.000,1.000) -- (3.000,1.000) -- (3.000,2.000) -- (2.000,2.000) -- cycle;
\node[elemlabel] at (1.500,1.500) {U};
\node[elemlabel] at (1.500,0.500) {L};
\draw[seledge] (1.000,1.000) -- (2.000,1.000);
\draw[hearrow, red!70!black] (1.220,1.140) -- (1.780,1.140);
\node[helabel, text=red!70!black] at (1.500,1.266) {$a$};
\draw[hearrow, blue!70!black] (1.780,0.860) -- (1.220,0.860);
\node[helabel, text=blue!70!black] at (1.500,0.734) {$b$};
\node[vreg, dot] at (0.000,0.000) {};
\node[vreg, dot] at (1.000,0.000) {};
\node[vreg, dot] at (2.000,0.000) {};
\node[vreg, dot] at (3.000,0.000) {};
\node[vreg, dot] at (0.000,1.000) {};
\node[vreg, minimum size=0.31cm] at (1.000,1.000) {$A$};
\node[vreg, minimum size=0.31cm] at (2.000,1.000) {$B$};
\node[vreg, dot] at (3.000,1.000) {};
\node[vreg, dot] at (0.000,2.000) {};
\node[vreg, dot] at (1.000,2.000) {};
\node[vreg, dot] at (2.000,2.000) {};
\node[vreg, dot] at (3.000,2.000) {};

\end{tikzpicture}
\caption{Antes: $A$ e $B$ internos.}\end{subfigure}\hfill
\begin{subfigure}[b]{0.28\linewidth}\centering\begin{tikzpicture}[scale=1.400, every node/.style={font=\scriptsize}]
\filldraw[cefill] (0.000,0.000) -- (1.000,0.000) -- (1.000,1.000) -- (0.000,1.000) -- cycle;
\filldraw[cefill] (1.000,0.000) -- (2.000,0.000) -- (2.000,1.000) -- (1.000,1.000) -- cycle;
\filldraw[cefill] (2.000,0.000) -- (3.000,0.000) -- (3.000,1.000) -- (2.000,1.000) -- cycle;
\filldraw[cefill] (0.000,1.000) -- (1.000,1.000) -- (1.000,2.000) -- (0.000,2.000) -- cycle;
\filldraw[cefill] (1.000,1.000) -- (2.000,1.000) -- (2.000,2.000) -- (1.000,2.000) -- cycle;
\filldraw[cefill] (2.000,1.000) -- (3.000,1.000) -- (3.000,2.000) -- (2.000,2.000) -- cycle;
\filldraw[iefill] (1.085,0.915) -- (1.915,0.915) -- (1.915,1.085) -- (1.085,1.085) -- cycle;
\draw[red, line width=1.1pt] (1.915,0.915) -- (1.915,1.085);
\draw[red, line width=1.1pt] (1.085,1.085) -- (1.085,0.915);
\node[elemlabel] at (1.500,1.500) {U};
\node[elemlabel] at (1.500,0.500) {L};
\draw[hearrow, red!70!black] (1.220,1.220) -- (1.780,1.220);
\node[helabel, text=red!70!black] at (1.276,1.352) {$a$};
\draw[hearrow, blue!70!black] (1.780,0.780) -- (1.220,0.780);
\node[helabel, text=blue!70!black] at (1.724,0.648) {$b$};
\draw[hearrow, black] (1.285,1.076) -- (1.285,0.924);
\node[helabel, text=black] at (0.865,1.200) {$h_3$};
\draw[hearrow, black] (1.715,0.924) -- (1.715,1.076);
\node[helabel, text=black] at (2.135,1.200) {$h_1$};
\node[vreg, dot] at (0.000,0.000) {};
\node[vreg, dot] at (1.000,0.000) {};
\node[vreg, dot] at (2.000,0.000) {};
\node[vreg, dot] at (3.000,0.000) {};
\node[vreg, dot] at (0.000,1.000) {};
\node[vcol, minimum size=0.31cm] at (1.000,1.000) {$A$};
\node[vcol, minimum size=0.31cm] at (2.000,1.000) {$B$};
\node[vreg, dot] at (3.000,1.000) {};
\node[vreg, dot] at (0.000,2.000) {};
\node[vreg, dot] at (1.000,2.000) {};
\node[vreg, dot] at (2.000,2.000) {};
\node[vreg, dot] at (3.000,2.000) {};

\end{tikzpicture}
\caption{Depois: elemento $[A,B,B,A]$.}\end{subfigure}\hfill
\begin{subfigure}[b]{0.2\linewidth}\centering\begin{tikzpicture}[scale=1.400, every node/.style={font=\scriptsize}]
\filldraw[cefill] (0.000,0.000) -- (1.000,0.000) -- (1.000,1.000) -- (0.000,1.000) -- cycle;
\filldraw[cefill] (1.000,0.000) -- (2.000,0.000) -- (2.000,1.000) -- (1.000,1.000) -- cycle;
\filldraw[cefill] (0.000,1.000) -- (1.000,1.000) -- (1.000,2.000) -- (0.000,2.000) -- cycle;
\filldraw[cefill] (1.000,1.000) -- (2.000,1.000) -- (2.000,2.000) -- (1.000,2.000) -- cycle;
\draw[seledge] (1.000,1.000) -- (1.000,2.000);
\draw[hearrow, red!70!black] (0.860,1.220) -- (0.860,1.780);
\node[helabel, text=red!70!black] at (0.734,1.500) {$a$};
\draw[hearrow, blue!70!black] (1.140,1.780) -- (1.140,1.220);
\node[helabel, text=blue!70!black] at (1.266,1.500) {$b$};
\node[vreg, dot] at (0.000,0.000) {};
\node[vreg, dot] at (1.000,0.000) {};
\node[vreg, dot] at (2.000,0.000) {};
\node[vreg, dot] at (0.000,1.000) {};
\node[vreg, minimum size=0.31cm] at (1.000,1.000) {$A$};
\node[vreg, dot] at (2.000,1.000) {};
\node[vreg, dot] at (0.000,2.000) {};
\node[vreg, minimum size=0.31cm] at (1.000,2.000) {$B$};
\node[vreg, dot] at (2.000,2.000) {};

\end{tikzpicture}
\caption{Antes: $B$ na borda.}\end{subfigure}\hfill
\begin{subfigure}[b]{0.2\linewidth}\centering\begin{tikzpicture}[scale=1.400, every node/.style={font=\scriptsize}]
\filldraw[cefill] (0.000,0.000) -- (1.000,0.000) -- (1.000,1.000) -- (0.000,1.000) -- cycle;
\filldraw[cefill] (1.000,0.000) -- (2.000,0.000) -- (2.000,1.000) -- (1.000,1.000) -- cycle;
\filldraw[cefill] (0.000,1.000) -- (1.000,1.000) -- (0.844,1.844) -- (0.000,2.000) -- cycle;
\filldraw[cefill] (1.000,1.000) -- (2.000,1.000) -- (2.000,2.000) -- (1.156,1.844) -- cycle;
\filldraw[iefill] (1.067,1.057) -- (1.156,1.844) -- (0.844,1.844) -- (0.933,1.057) -- cycle;
\draw[red, line width=1.1pt] (0.933,1.057) -- (1.067,1.057);
\draw[hearrow, red!70!black] (0.826,1.198) -- (0.739,1.670);
\node[helabel, text=red!70!black] at (0.657,1.445) {$a$};
\draw[hearrow, blue!70!black] (1.261,1.670) -- (1.174,1.198);
\node[helabel, text=blue!70!black] at (1.343,1.445) {$b$};
\node[vreg, dot] at (0.000,0.000) {};
\node[vreg, dot] at (1.000,0.000) {};
\node[vreg, dot] at (2.000,0.000) {};
\node[vreg, dot] at (0.000,1.000) {};
\node[vcol, minimum size=0.31cm] at (1.000,1.000) {$A$};
\node[vreg, dot] at (2.000,1.000) {};
\node[vreg, dot] at (0.000,2.000) {};
\node[vie, minimum size=0.31cm] at (0.844,1.844) {$B$};
\node[vreg, dot] at (2.000,2.000) {};
\node[vie, minimum size=0.31cm] at (1.156,1.844) {$B_{1}$};

\end{tikzpicture}
\caption{Depois.}\end{subfigure}
\caption{Operador Canônico. (a--b) Com os dois vértices internos, o elemento de interface tem as duas extremidades
colapsadas ($h_3$ em $A$ e $h_1$ em $B$, em vermelho) e nenhum vértice é criado: geometricamente ele é só a aresta $AB$.
(c--d) Com $B$ na borda, $B$ é duplicado em $B_1$; o operador em $B$ recebe $b$, e os elementos do lado de $b$ (à
direita) passam a usar a cópia; o lado curto em $B$ é uma aresta de borda. Em $A$ (interno) a extremidade continua colapsada.}
\label{fig:canonico}
\end{figure}

\subsection{Split Element}
\label{sec:split}

\textbf{Quando:} $v$ tem uma aresta colapsada $c$ (pertence a exatamente um elemento de interface $I$, criado pelo
Canônico) \emph{e} $I$ é vizinho de um dos dois elementos que contêm a aresta nova. Esse elemento comum é o
\emph{elemento compartilhado} $S$.

\textbf{Como o código testa} (\code{retrieveAvailableVertex}, \cref{alg:retrieve}). Seja $he$ a half-edge nova que parte
de $v$ e $b=\Op[he]$. Há duas posições possíveis para $S$:
\begin{itemize}
\item \textbf{hipótese $x_1$}: $S$ é o elemento de $he$. Então a aresta de $S$ que chega em $v$, $\prv(he)$, é vizinha
de $I$, e sua oposta $x_1=\Op[\prv(he)]$ é o lado longo de $I$ que \emph{sai} de $v$ logo depois de $c$, isto é,
$\prv(x_1)=c$;
\item \textbf{hipótese $x_3$}: $S$ é o elemento de $b$. Então a aresta de $S$ que sai de $v$, $\nxt(b)$, é vizinha de
$I$, e $x_3=\Op[\nxt(b)]$ é o lado longo de $I$ que \emph{chega} em $v$ logo antes de $c$, isto é, $\nxt(x_3)=c$.
\end{itemize}
Qual hipótese ocorre não depende de qual half-edge da aresta o usuário escolheu, e sim da posição da aresta nova em
relação a $I$: em $S$, se a aresta nova vem depois da aresta de $I$ (no sentido anti-horário de $S$) vale $x_1$; se vem
antes, vale $x_3$. A \cref{fig:hipoteses} mostra, em detalhe, as consultas feitas em cada caso e também o caso em que
nenhuma das duas vale; a \cref{fig:split} mostra o resultado numa malha.

% Base comum das três figuras: vértice v na origem, elemento de interface I à esquerda (colapsado em v) e quatro
% elementos contínuos em volta. \hipS{x}{y} rotula o elemento compartilhado.
\newcommand{\hipbase}[1]{%
  \filldraw[cefill] (-2.8,0.6) -- (0,0) -- (0,2.4) -- (-2.8,2.4) -- cycle;
  \filldraw[cefill] (-2.8,-2.4) -- (0,-2.4) -- (0,0) -- (-2.8,-0.6) -- cycle;
  \filldraw[cefill] (0,0) -- (2.6,0) -- (2.6,2.4) -- (0,2.4) -- cycle;
  \filldraw[cefill] (0,-2.4) -- (2.6,-2.4) -- (2.6,0) -- (0,0) -- cycle;
  \filldraw[iefill] (-2.8,-0.6) -- (0,0) -- (-2.8,0.6) -- cycle;
  \draw[white, line width=1.2pt] (-2.8,-0.55) -- (-2.8,0.55);
  \draw[orange!55!black, dashed, line width=0.45pt] (-2.8,-0.6) -- (-2.8,0.6);
  \node[font=\footnotesize\bfseries, text=orange!40!black] at (-2.35,0) {$I$};
  #1
  \node[vcol, minimum size=0.32cm] at (0,0) {};
  \node[font=\scriptsize] at (0.28,-0.22) {$v$};
}
\newcommand{\hipc}[2]{% rótulo da aresta colapsada c, com uma linha até v
  \draw[red!70!black, line width=0.4pt] (0.05,0.03*#2) -- (#1,0.5*#2);
  \node[font=\scriptsize, text=red!70!black, anchor=west, align=left] at (#1,0.5*#2) {$c$ (colapsada,\\ de $v$ para $v$)};
}

\begin{figure}[ht]
\centering
\begin{subfigure}[b]{0.325\linewidth}\centering
\begin{tikzpicture}[scale=0.88, every node/.style={font=\scriptsize}]
  \hipbase{\draw[curedge] (0,0) -- (0,2.4);}
  \node[elemlabel] at (-1.5,1.65) {S};
  \draw[hearrow, red!70!black, line width=0.8pt] (-0.17,0.45) -- (-0.17,2.0) node[midway, left=0pt] {$he$};
  \draw[hearrow, blue!70!black] (0.17,2.0) -- (0.17,0.45) node[midway, right=0pt] {$b$};
  \draw[hearrow, violet, line width=0.8pt] (-2.4,0.68) -- (-0.5,0.28) node[pos=0.45, above=1pt, sloped] {$\prv(he)$};
  \draw[hearrow, green!45!black, line width=0.8pt] (-0.9,0.07) -- (-2.5,0.42);
  \node[text=green!45!black] at (-1.75,0.03) {$x_1$};
  \hipc{0.75}{-1}
\end{tikzpicture}
\caption{Hipótese $x_1$: $S$ é o elemento de $he$.\\ $x_1=\Op[\prv(he)]$ e $\prv(x_1)=c$.}
\end{subfigure}\hfill
\begin{subfigure}[b]{0.325\linewidth}\centering
\begin{tikzpicture}[scale=0.88, every node/.style={font=\scriptsize}]
  \hipbase{\draw[curedge] (0,0) -- (0,-2.4);}
  \node[elemlabel] at (-1.5,-1.65) {S};
  \draw[hearrow, red!70!black, line width=0.8pt] (0.17,-0.45) -- (0.17,-2.0) node[midway, right=0pt] {$he$};
  \draw[hearrow, blue!70!black] (-0.17,-2.0) -- (-0.17,-0.45) node[midway, left=0pt] {$b$};
  \draw[hearrow, violet, line width=0.8pt] (-0.5,-0.28) -- (-2.4,-0.68) node[pos=0.55, below=1pt, sloped] {$\nxt(b)$};
  \draw[hearrow, green!45!black, line width=0.8pt] (-2.5,-0.42) -- (-0.9,-0.07);
  \node[text=green!45!black] at (-1.75,-0.03) {$x_3$};
  \hipc{0.75}{1}
\end{tikzpicture}
\caption{Hipótese $x_3$: $S$ é o elemento de $b$.\\ $x_3=\Op[\nxt(b)]$ e $\nxt(x_3)=c$.}
\end{subfigure}\hfill
\begin{subfigure}[b]{0.325\linewidth}\centering
\begin{tikzpicture}[scale=0.88, every node/.style={font=\scriptsize}]
  \hipbase{\draw[curedge] (0,0) -- (2.6,0);}
  \draw[hearrow, red!70!black, line width=0.8pt] (0.45,0.17) -- (2.2,0.17) node[midway, above=0pt] {$he$};
  \draw[hearrow, blue!70!black] (2.2,-0.17) -- (0.45,-0.17) node[midway, below=0pt] {$b$};
  \draw[hearrow, violet, line width=0.8pt] (0.17,2.0) -- (0.17,0.45) node[pos=0.3, right=0pt] {$\prv(he)$};
  \draw[hearrow, black!50, dashed] (-0.17,0.45) -- (-0.17,2.0) node[pos=0.7, left=0pt, text=black!60] {$\Op[\prv(he)]$};
  \draw[hearrow, violet, line width=0.8pt] (0.17,-0.45) -- (0.17,-2.0) node[pos=0.7, right=0pt] {$\nxt(b)$};
  \draw[hearrow, black!50, dashed] (-0.17,-2.0) -- (-0.17,-0.45) node[pos=0.3, left=0pt, text=black!60] {$\Op[\nxt(b)]$};
\end{tikzpicture}
\caption{Nenhuma: os opostos caem fora de $I$.\\ $\Rightarrow$ Expand Edge.}
\end{subfigure}
\caption{As consultas de \code{retrieveAvailableVertex} em torno de um vértice $v$ que tem uma aresta colapsada $c$ (o
elemento de interface $I$, à esquerda, termina em $v$; a aresta nova está em vermelho grosso). (a) A aresta nova sobe:
o elemento de $he$ tem uma aresta, $\prv(he)$, colada ao lado longo de $I$ que sai de $v$ logo depois de $c$. (b) A
aresta nova desce: agora é o elemento de $b$ que encosta em $I$, pela aresta $\nxt(b)$, colada ao lado longo de $I$ que
chega em $v$ logo antes de $c$. (c) A aresta nova continua a trinca: os dois testes levam a elementos contínuos, e o
operador escolhido é o Expand Edge. Em (a) e (b), o vértice novo vai para $S$.}
\label{fig:hipoteses}
\end{figure}

\textbf{O que faz} (\cref{alg:split}): cria um vértice $n$ e o coloca ``dentro'' de $S$: só dois cantos mudam, o canto de
$S$ em $v$ e o canto de $I$ adjacente a $S$. A antiga aresta colapsada $c$ deixa de ter comprimento zero (passa a ligar
$v$ e $n$) e vira a fronteira entre $I$ e o novo elemento, sendo ligada a \code{elementHE}. O custo é um número
constante de passos ($O(\log n_c)$ na IBHM, por causa das consultas $\nxt$/$\prv$).

\begin{algorithm}[ht]
\caption{\code{retrieveAvailableVertex}$(he)$ e \code{splitElement}$(he, \mathit{av}, \mathit{sh}, x)$}
\label{alg:retrieve}\label{alg:split}
\begin{algorithmic}[1]
\Function{retrieveAvailableVertex}{$he$}
  \State $v\gets\V[he]$; \quad $b\gets\Op[he]$; \quad $c\gets$ \code{\_collapsed}$[v]$
  \State $x_1\gets\Op[\prv(he)]$
  \If{$x_1\ne\BORDER$ \textbf{e} $\prv(x_1)=c$} \Return $(\mathit{av}=x_1,\ \mathit{sh}=he)$ \Comment{hipótese $x_1$}
  \EndIf
  \State $x_3\gets\Op[\nxt(b)]$
  \If{$x_3\ne\BORDER$ \textbf{e} $\nxt(x_3)=c$} \Return $(\mathit{av}=c,\ \mathit{sh}=\nxt(b))$ \Comment{hipótese $x_3$}
  \EndIf
  \State \Return $(\BORDER,\BORDER)$ \Comment{nenhum elemento compartilhado $\Rightarrow$ Expand Edge}
\EndFunction
\Statex
\Procedure{splitElement}{$he, \mathit{av}, \mathit{sh}, x$}
  \State $v\gets\V[he]$; \quad $n\gets$ \code{duplicatePoint}$(v)$
  \State $\V[\mathit{av}]\gets n$ \Comment{canto do elemento de interface antigo}
  \State $\V[\mathit{sh}]\gets n$ \Comment{canto do elemento compartilhado}
  \If{$he=\mathit{sh}$} \Comment{hipótese $x_1$}
    \State $\V[x]\gets n$; \quad $\V[\nxt(x)]\gets v$; \quad $\Op[x]\leftrightarrow\prv(\mathit{av})$ \Comment{$\prv(\mathit{av})=c$}
  \Else \Comment{hipótese $x_3$}
    \State $\V[x]\gets v$; \quad $\V[\nxt(x)]\gets n$; \quad $\Op[x]\leftrightarrow \mathit{av}$ \Comment{$\mathit{av}=c$}
  \EndIf
  \State marca $v$ e $n$ como vértices de EI; remove $v$ de \code{\_collapsed}
\EndProcedure
\end{algorithmic}
\end{algorithm}

\begin{figure}[ht]
\centering
\begin{subfigure}[b]{0.31\linewidth}\centering\begin{tikzpicture}[scale=2.000, every node/.style={font=\scriptsize}]
\clip (0.450,0.550) rectangle (2.800,3.450);
\filldraw[cefill] (0.000,0.000) -- (1.000,0.000) -- (1.000,1.000) -- (0.000,1.000) -- cycle;
\filldraw[cefill] (1.000,0.000) -- (2.000,0.000) -- (2.000,1.000) -- (1.000,1.000) -- cycle;
\filldraw[cefill] (2.000,0.000) -- (3.000,0.000) -- (3.000,1.000) -- (2.000,1.000) -- cycle;
\filldraw[cefill] (0.000,1.000) -- (1.000,1.000) -- (1.000,2.000) -- (0.000,2.000) -- cycle;
\filldraw[cefill] (1.000,1.000) -- (2.000,1.000) -- (2.000,2.000) -- (1.000,2.000) -- cycle;
\filldraw[cefill] (2.000,1.000) -- (3.000,1.000) -- (3.000,2.000) -- (2.000,2.000) -- cycle;
\filldraw[cefill] (0.000,2.000) -- (1.000,2.000) -- (1.000,3.000) -- (0.000,3.000) -- cycle;
\filldraw[cefill] (1.000,2.000) -- (2.000,2.000) -- (2.000,3.000) -- (1.000,3.000) -- cycle;
\filldraw[cefill] (2.000,2.000) -- (3.000,2.000) -- (3.000,3.000) -- (2.000,3.000) -- cycle;
\filldraw[cefill] (0.000,3.000) -- (1.000,3.000) -- (1.000,4.000) -- (0.000,4.000) -- cycle;
\filldraw[cefill] (1.000,3.000) -- (2.000,3.000) -- (2.000,4.000) -- (1.000,4.000) -- cycle;
\filldraw[cefill] (2.000,3.000) -- (3.000,3.000) -- (3.000,4.000) -- (2.000,4.000) -- cycle;
\filldraw[iefill] (1.062,1.938) -- (1.938,1.938) -- (1.938,2.062) -- (1.062,2.062) -- cycle;
\draw[red, line width=1.1pt] (1.938,1.938) -- (1.938,2.062);
\draw[red, line width=1.1pt] (1.062,2.062) -- (1.062,1.938);
\node[elemlabel] at (1.500,2.500) {S};
\draw[seledge] (2.000,2.000) -- (2.000,3.000);
\draw[hearrow, red!70!black] (1.860,2.220) -- (1.860,2.780);
\node[helabel, text=red!70!black] at (1.692,2.500) {$he$};
\draw[hearrow, blue!70!black] (2.140,2.780) -- (2.140,2.220);
\node[helabel, text=blue!70!black] at (2.266,2.500) {$b$};
\draw[hearrow, violet] (1.220,2.140) -- (1.780,2.140);
\node[helabel, text=violet] at (1.500,2.364) {$\mathrm{prev}(he)$};
\draw[hearrow, green!45!black] (1.745,2.034) -- (1.255,2.034);
\node[helabel, text=green!45!black] at (1.500,1.961) {$x_1$};
\node[vreg, dot] at (0.000,0.000) {};
\node[vreg, dot] at (1.000,0.000) {};
\node[vreg, dot] at (2.000,0.000) {};
\node[vreg, dot] at (3.000,0.000) {};
\node[vreg, dot] at (0.000,1.000) {};
\node[vreg, dot] at (1.000,1.000) {};
\node[vreg, minimum size=0.29cm] at (2.000,1.000) {$D$};
\node[vreg, dot] at (3.000,1.000) {};
\node[vreg, dot] at (0.000,2.000) {};
\node[vcol, minimum size=0.29cm] at (1.000,2.000) {$A$};
\node[vcol, minimum size=0.29cm] at (2.000,2.000) {$B$};
\node[vreg, dot] at (3.000,2.000) {};
\node[vreg, dot] at (0.000,3.000) {};
\node[vreg, dot] at (1.000,3.000) {};
\node[vreg, minimum size=0.29cm] at (2.000,3.000) {$C$};
\node[vreg, dot] at (3.000,3.000) {};
\node[vreg, dot] at (0.000,4.000) {};
\node[vreg, dot] at (1.000,4.000) {};
\node[vreg, dot] at (2.000,4.000) {};
\node[vreg, dot] at (3.000,4.000) {};

\end{tikzpicture}
\caption{$x_1$: antes.}\end{subfigure}\hspace{0.06\linewidth}
\begin{subfigure}[b]{0.31\linewidth}\centering\begin{tikzpicture}[scale=2.000, every node/.style={font=\scriptsize}]
\clip (0.450,0.550) rectangle (2.800,3.450);
\filldraw[cefill] (0.000,0.000) -- (1.000,0.000) -- (1.000,1.000) -- (0.000,1.000) -- cycle;
\filldraw[cefill] (1.000,0.000) -- (2.000,0.000) -- (2.000,1.000) -- (1.000,1.000) -- cycle;
\filldraw[cefill] (2.000,0.000) -- (3.000,0.000) -- (3.000,1.000) -- (2.000,1.000) -- cycle;
\filldraw[cefill] (0.000,1.000) -- (1.000,1.000) -- (1.000,2.000) -- (0.000,2.000) -- cycle;
\filldraw[cefill] (1.000,1.000) -- (2.000,1.000) -- (2.156,1.844) -- (1.000,2.000) -- cycle;
\filldraw[cefill] (2.000,1.000) -- (3.000,1.000) -- (3.000,2.000) -- (2.156,1.844) -- cycle;
\filldraw[cefill] (0.000,2.000) -- (1.000,2.000) -- (1.000,3.000) -- (0.000,3.000) -- cycle;
\filldraw[cefill] (1.000,2.000) -- (1.844,2.156) -- (2.000,3.000) -- (1.000,3.000) -- cycle;
\filldraw[cefill] (2.156,1.844) -- (3.000,2.000) -- (3.000,3.000) -- (2.000,3.000) -- cycle;
\filldraw[cefill] (0.000,3.000) -- (1.000,3.000) -- (1.000,4.000) -- (0.000,4.000) -- cycle;
\filldraw[cefill] (1.000,3.000) -- (2.000,3.000) -- (2.000,4.000) -- (1.000,4.000) -- cycle;
\filldraw[cefill] (2.000,3.000) -- (3.000,3.000) -- (3.000,4.000) -- (2.000,4.000) -- cycle;
\filldraw[iefill] (1.062,1.938) -- (2.156,1.844) -- (1.844,2.156) -- (1.057,2.067) -- cycle;
\filldraw[iefill] (2.156,1.844) -- (2.062,2.938) -- (1.933,2.943) -- (1.844,2.156) -- cycle;
\draw[red, line width=1.1pt] (1.057,2.067) -- (1.062,1.938);
\draw[red, line width=1.1pt] (2.062,2.938) -- (1.933,2.943);
\node[elemlabel] at (1.461,2.539) {S};
\node[elemlabel] at (1.999,2.470) {};
\node[vreg, dot] at (0.000,0.000) {};
\node[vreg, dot] at (1.000,0.000) {};
\node[vreg, dot] at (2.000,0.000) {};
\node[vreg, dot] at (3.000,0.000) {};
\node[vreg, dot] at (0.000,1.000) {};
\node[vreg, dot] at (1.000,1.000) {};
\node[vreg, minimum size=0.29cm] at (2.000,1.000) {$D$};
\node[vreg, dot] at (3.000,1.000) {};
\node[vreg, dot] at (0.000,2.000) {};
\node[vcol, minimum size=0.29cm] at (1.000,2.000) {$A$};
\node[vie, minimum size=0.29cm] at (2.156,1.844) {$B$};
\node[vreg, dot] at (3.000,2.000) {};
\node[vreg, dot] at (0.000,3.000) {};
\node[vreg, dot] at (1.000,3.000) {};
\node[vcol, minimum size=0.29cm] at (2.000,3.000) {$C$};
\node[vreg, dot] at (3.000,3.000) {};
\node[vreg, dot] at (0.000,4.000) {};
\node[vreg, dot] at (1.000,4.000) {};
\node[vreg, dot] at (2.000,4.000) {};
\node[vreg, dot] at (3.000,4.000) {};
\node[vie, minimum size=0.29cm] at (1.844,2.156) {$B_{1}$};

\end{tikzpicture}
\caption{$x_1$: depois.}\end{subfigure}\\[6pt]
\begin{subfigure}[b]{0.31\linewidth}\centering\begin{tikzpicture}[scale=2.000, every node/.style={font=\scriptsize}]
\clip (0.450,0.550) rectangle (2.800,3.450);
\filldraw[cefill] (0.000,0.000) -- (1.000,0.000) -- (1.000,1.000) -- (0.000,1.000) -- cycle;
\filldraw[cefill] (1.000,0.000) -- (2.000,0.000) -- (2.000,1.000) -- (1.000,1.000) -- cycle;
\filldraw[cefill] (2.000,0.000) -- (3.000,0.000) -- (3.000,1.000) -- (2.000,1.000) -- cycle;
\filldraw[cefill] (0.000,1.000) -- (1.000,1.000) -- (1.000,2.000) -- (0.000,2.000) -- cycle;
\filldraw[cefill] (1.000,1.000) -- (2.000,1.000) -- (2.000,2.000) -- (1.000,2.000) -- cycle;
\filldraw[cefill] (2.000,1.000) -- (3.000,1.000) -- (3.000,2.000) -- (2.000,2.000) -- cycle;
\filldraw[cefill] (0.000,2.000) -- (1.000,2.000) -- (1.000,3.000) -- (0.000,3.000) -- cycle;
\filldraw[cefill] (1.000,2.000) -- (2.000,2.000) -- (2.000,3.000) -- (1.000,3.000) -- cycle;
\filldraw[cefill] (2.000,2.000) -- (3.000,2.000) -- (3.000,3.000) -- (2.000,3.000) -- cycle;
\filldraw[cefill] (0.000,3.000) -- (1.000,3.000) -- (1.000,4.000) -- (0.000,4.000) -- cycle;
\filldraw[cefill] (1.000,3.000) -- (2.000,3.000) -- (2.000,4.000) -- (1.000,4.000) -- cycle;
\filldraw[cefill] (2.000,3.000) -- (3.000,3.000) -- (3.000,4.000) -- (2.000,4.000) -- cycle;
\filldraw[iefill] (1.062,1.938) -- (1.938,1.938) -- (1.938,2.062) -- (1.062,2.062) -- cycle;
\draw[red, line width=1.1pt] (1.938,1.938) -- (1.938,2.062);
\draw[red, line width=1.1pt] (1.062,2.062) -- (1.062,1.938);
\node[elemlabel] at (1.500,1.500) {S};
\draw[seledge] (2.000,2.000) -- (2.000,1.000);
\draw[hearrow, red!70!black] (2.140,1.780) -- (2.140,1.220);
\node[helabel, text=red!70!black] at (2.308,1.500) {$he$};
\draw[hearrow, blue!70!black] (1.860,1.220) -- (1.860,1.780);
\node[helabel, text=blue!70!black] at (1.734,1.500) {$b$};
\draw[hearrow, violet] (1.780,1.860) -- (1.220,1.860);
\node[helabel, text=violet] at (1.500,1.636) {$\mathrm{next}(b)$};
\draw[hearrow, green!45!black] (1.255,1.966) -- (1.745,1.966);
\node[helabel, text=green!45!black] at (1.500,2.039) {$x_3$};
\node[vreg, dot] at (0.000,0.000) {};
\node[vreg, dot] at (1.000,0.000) {};
\node[vreg, dot] at (2.000,0.000) {};
\node[vreg, dot] at (3.000,0.000) {};
\node[vreg, dot] at (0.000,1.000) {};
\node[vreg, dot] at (1.000,1.000) {};
\node[vreg, minimum size=0.29cm] at (2.000,1.000) {$D$};
\node[vreg, dot] at (3.000,1.000) {};
\node[vreg, dot] at (0.000,2.000) {};
\node[vcol, minimum size=0.29cm] at (1.000,2.000) {$A$};
\node[vcol, minimum size=0.29cm] at (2.000,2.000) {$B$};
\node[vreg, dot] at (3.000,2.000) {};
\node[vreg, dot] at (0.000,3.000) {};
\node[vreg, dot] at (1.000,3.000) {};
\node[vreg, minimum size=0.29cm] at (2.000,3.000) {$C$};
\node[vreg, dot] at (3.000,3.000) {};
\node[vreg, dot] at (0.000,4.000) {};
\node[vreg, dot] at (1.000,4.000) {};
\node[vreg, dot] at (2.000,4.000) {};
\node[vreg, dot] at (3.000,4.000) {};

\end{tikzpicture}
\caption{$x_3$: antes.}\end{subfigure}\hspace{0.06\linewidth}
\begin{subfigure}[b]{0.31\linewidth}\centering\begin{tikzpicture}[scale=2.000, every node/.style={font=\scriptsize}]
\clip (0.450,0.550) rectangle (2.800,3.450);
\filldraw[cefill] (0.000,0.000) -- (1.000,0.000) -- (1.000,1.000) -- (0.000,1.000) -- cycle;
\filldraw[cefill] (1.000,0.000) -- (2.000,0.000) -- (2.000,1.000) -- (1.000,1.000) -- cycle;
\filldraw[cefill] (2.000,0.000) -- (3.000,0.000) -- (3.000,1.000) -- (2.000,1.000) -- cycle;
\filldraw[cefill] (0.000,1.000) -- (1.000,1.000) -- (1.000,2.000) -- (0.000,2.000) -- cycle;
\filldraw[cefill] (1.000,1.000) -- (2.000,1.000) -- (1.844,1.844) -- (1.000,2.000) -- cycle;
\filldraw[cefill] (2.000,1.000) -- (3.000,1.000) -- (3.000,2.000) -- (2.156,2.156) -- cycle;
\filldraw[cefill] (0.000,2.000) -- (1.000,2.000) -- (1.000,3.000) -- (0.000,3.000) -- cycle;
\filldraw[cefill] (1.000,2.000) -- (2.156,2.156) -- (2.000,3.000) -- (1.000,3.000) -- cycle;
\filldraw[cefill] (2.156,2.156) -- (3.000,2.000) -- (3.000,3.000) -- (2.000,3.000) -- cycle;
\filldraw[cefill] (0.000,3.000) -- (1.000,3.000) -- (1.000,4.000) -- (0.000,4.000) -- cycle;
\filldraw[cefill] (1.000,3.000) -- (2.000,3.000) -- (2.000,4.000) -- (1.000,4.000) -- cycle;
\filldraw[cefill] (2.000,3.000) -- (3.000,3.000) -- (3.000,4.000) -- (2.000,4.000) -- cycle;
\filldraw[iefill] (1.057,1.933) -- (1.844,1.844) -- (2.156,2.156) -- (1.062,2.062) -- cycle;
\filldraw[iefill] (1.844,1.844) -- (1.933,1.057) -- (2.062,1.062) -- (2.156,2.156) -- cycle;
\draw[red, line width=1.1pt] (1.062,2.062) -- (1.057,1.933);
\draw[red, line width=1.1pt] (1.933,1.057) -- (2.062,1.062);
\node[elemlabel] at (1.461,1.461) {S};
\node[elemlabel] at (1.999,1.530) {};
\node[vreg, dot] at (0.000,0.000) {};
\node[vreg, dot] at (1.000,0.000) {};
\node[vreg, dot] at (2.000,0.000) {};
\node[vreg, dot] at (3.000,0.000) {};
\node[vreg, dot] at (0.000,1.000) {};
\node[vreg, dot] at (1.000,1.000) {};
\node[vcol, minimum size=0.29cm] at (2.000,1.000) {$D$};
\node[vreg, dot] at (3.000,1.000) {};
\node[vreg, dot] at (0.000,2.000) {};
\node[vcol, minimum size=0.29cm] at (1.000,2.000) {$A$};
\node[vie, minimum size=0.29cm] at (2.156,2.156) {$B$};
\node[vreg, dot] at (3.000,2.000) {};
\node[vreg, dot] at (0.000,3.000) {};
\node[vreg, dot] at (1.000,3.000) {};
\node[vreg, minimum size=0.29cm] at (2.000,3.000) {$C$};
\node[vreg, dot] at (3.000,3.000) {};
\node[vreg, dot] at (0.000,4.000) {};
\node[vreg, dot] at (1.000,4.000) {};
\node[vreg, dot] at (2.000,4.000) {};
\node[vreg, dot] at (3.000,4.000) {};
\node[vie, minimum size=0.29cm] at (1.844,1.844) {$B_{1}$};

\end{tikzpicture}
\caption{$x_3$: depois.}\end{subfigure}
\caption{Split Element em $B$. Nos dois casos já existe um elemento de interface em $AB$, colapsado em $A$ e em $B$.
(a--b) A aresta nova $BC$ sobe: o elemento compartilhado $S$ (acima de $AB$) contém $he$, e o lado de $I$ que sai de $B$
é $x_1$. (c--d) A aresta nova $BD$ desce: agora $S$ (abaixo de $AB$) contém $b$, e o lado de $I$ que chega em $B$ é
$x_3$. Em ambos, só $S$ e um canto de $I$ passam a usar a nova cópia $B_1$, e a antiga aresta colapsada em $B$ vira o
lado comum entre os dois elementos de interface. Note que $C$ e $D$, internos, ficam colapsados.}
\label{fig:split}
\end{figure}

\subsection{Expand Edge}
\label{sec:expand}

\textbf{Quando:} $v$ tem uma aresta colapsada $c$, mas nenhuma das duas hipóteses vale: o elemento de interface $I$ não
toca os elementos da aresta nova. É o caso típico de uma trinca que continua ``em frente'' através de $v$
(\cref{fig:expand}).

\textbf{O que faz} (\cref{alg:expand}): a estrela de $v$ precisa ser dividida em dois setores, delimitados pela aresta
nova de um lado e por $I$ do outro. O operador cria $n$ e gira no sentido anti-horário \emph{a partir da aresta
colapsada} $c$: primeiro troca o canto de $c$ em $I$, depois sai de $I$ pelo seu lado longo $\prv(c)$ e percorre os
elementos daquele setor até alcançar $he$, que \emph{não} é alterado. Assim, a parada natural da rotação é justamente
cruzar a aresta nova. A aresta $c$ deixa de ser colapsada e é ligada a \code{elementHE}.

\begin{algorithm}[ht]
\caption{\code{expandEdge}$(he, x)$}
\label{alg:expand}
\begin{algorithmic}[1]
\State $v\gets\V[he]$; \quad $n\gets$ \code{duplicatePoint}$(v)$; \quad $c\gets$ \code{\_collapsed}$[v]$
\State \code{reindexElementsCCW}$(c, he, n)$ \Comment{de $c$ até $he$ (exclusive)}
\State $\Op[c]\leftrightarrow x$
\State $\V[x]\gets v$; \quad $\V[\nxt(x)]\gets n$
\State marca $v$ e $n$ como vértices de EI; remove $v$ de \code{\_collapsed}
\end{algorithmic}
\end{algorithm}

\begin{figure}[ht]
\centering
\begin{subfigure}[b]{0.32\linewidth}\centering\resizebox{\linewidth}{!}{\begin{tikzpicture}[scale=1.700, every node/.style={font=\scriptsize}]
\clip (0.450,0.300) rectangle (3.550,1.700);
\filldraw[cefill] (0.000,0.000) -- (1.000,0.000) -- (1.000,1.000) -- (0.000,1.000) -- cycle;
\filldraw[cefill] (1.000,0.000) -- (2.000,0.000) -- (2.000,1.000) -- (1.000,1.000) -- cycle;
\filldraw[cefill] (2.000,0.000) -- (3.000,0.000) -- (3.000,1.000) -- (2.000,1.000) -- cycle;
\filldraw[cefill] (3.000,0.000) -- (4.000,0.000) -- (4.000,1.000) -- (3.000,1.000) -- cycle;
\filldraw[cefill] (0.000,1.000) -- (1.000,1.000) -- (1.000,2.000) -- (0.000,2.000) -- cycle;
\filldraw[cefill] (1.000,1.000) -- (2.000,1.000) -- (2.000,2.000) -- (1.000,2.000) -- cycle;
\filldraw[cefill] (2.000,1.000) -- (3.000,1.000) -- (3.000,2.000) -- (2.000,2.000) -- cycle;
\filldraw[cefill] (3.000,1.000) -- (4.000,1.000) -- (4.000,2.000) -- (3.000,2.000) -- cycle;
\filldraw[iefill] (1.057,0.943) -- (1.943,0.943) -- (1.943,1.057) -- (1.057,1.057) -- cycle;
\draw[red, line width=1.1pt] (1.943,0.943) -- (1.943,1.057);
\draw[red, line width=1.1pt] (1.057,1.057) -- (1.057,0.943);
\draw[seledge] (2.000,1.000) -- (3.000,1.000);
\draw[hearrow, red!70!black] (2.220,1.140) -- (2.780,1.140);
\node[helabel, text=red!70!black] at (2.500,1.266) {$he$};
\draw[hearrow, orange!80!black] (1.803,0.968) -- (1.803,1.032);
\node[helabel, text=orange!80!black] at (1.677,1.000) {$c$};
\draw[hearrow, blue!70!black] (2.780,0.860) -- (2.220,0.860);
\node[helabel, text=blue!70!black] at (2.500,0.734) {$b$};
\node[vreg, dot] at (0.000,0.000) {};
\node[vreg, dot] at (1.000,0.000) {};
\node[vreg, dot] at (2.000,0.000) {};
\node[vreg, dot] at (3.000,0.000) {};
\node[vreg, dot] at (4.000,0.000) {};
\node[vreg, dot] at (0.000,1.000) {};
\node[vcol, minimum size=0.29cm] at (1.000,1.000) {$A$};
\node[vcol, minimum size=0.29cm] at (2.000,1.000) {$B$};
\node[vreg, minimum size=0.29cm] at (3.000,1.000) {$C$};
\node[vreg, dot] at (4.000,1.000) {};
\node[vreg, dot] at (0.000,2.000) {};
\node[vreg, dot] at (1.000,2.000) {};
\node[vreg, dot] at (2.000,2.000) {};
\node[vreg, dot] at (3.000,2.000) {};
\node[vreg, dot] at (4.000,2.000) {};

\end{tikzpicture}
}\caption{Antes.}\end{subfigure}\hfill
\begin{subfigure}[b]{0.32\linewidth}\centering\resizebox{\linewidth}{!}{\begin{tikzpicture}[scale=1.700, every node/.style={font=\scriptsize}]
\clip (0.450,0.300) rectangle (3.550,1.700);
\filldraw[cefill] (0.000,0.000) -- (1.000,0.000) -- (1.000,1.000) -- (0.000,1.000) -- cycle;
\filldraw[cefill] (1.000,0.000) -- (2.000,0.000) -- (2.000,1.000) -- (1.000,1.000) -- cycle;
\filldraw[cefill] (2.000,0.000) -- (3.000,0.000) -- (3.000,1.000) -- (2.000,1.000) -- cycle;
\filldraw[cefill] (3.000,0.000) -- (4.000,0.000) -- (4.000,1.000) -- (3.000,1.000) -- cycle;
\filldraw[cefill] (0.000,1.000) -- (1.000,1.000) -- (1.000,2.000) -- (0.000,2.000) -- cycle;
\filldraw[cefill] (1.000,1.000) -- (2.000,1.000) -- (2.000,2.000) -- (1.000,2.000) -- cycle;
\filldraw[cefill] (2.000,1.000) -- (3.000,1.000) -- (3.000,2.000) -- (2.000,2.000) -- cycle;
\filldraw[cefill] (3.000,1.000) -- (4.000,1.000) -- (4.000,2.000) -- (3.000,2.000) -- cycle;
\filldraw[iefill] (1.057,0.943) -- (1.943,0.943) -- (1.943,1.057) -- (1.057,1.057) -- cycle;
\draw[red, line width=1.1pt] (1.943,0.943) -- (1.943,1.057);
\draw[red, line width=1.1pt] (1.057,1.057) -- (1.057,0.943);
\draw[seledge] (2.000,1.000) -- (3.000,1.000);
\draw[hearrow, orange!80!black] (1.803,0.968) -- (1.803,1.032);
\node[helabel, text=orange!80!black] at (1.583,1.200) {$1$};
\draw[hearrow, orange!80!black] (1.780,0.860) -- (1.220,0.860);
\node[helabel, text=orange!80!black] at (1.500,0.664) {$2$};
\draw[hearrow, orange!80!black] (2.140,0.780) -- (2.140,0.220);
\node[helabel, text=orange!80!black] at (2.336,0.500) {$3$};
\draw[hearrow, red!70!black] (2.220,1.140) -- (2.780,1.140);
\node[helabel, text=red!70!black] at (2.500,1.322) {$he=\mathrm{stop}$};
\node[vreg, dot] at (0.000,0.000) {};
\node[vreg, dot] at (1.000,0.000) {};
\node[vreg, dot] at (2.000,0.000) {};
\node[vreg, dot] at (3.000,0.000) {};
\node[vreg, dot] at (4.000,0.000) {};
\node[vreg, dot] at (0.000,1.000) {};
\node[vcol, minimum size=0.29cm] at (1.000,1.000) {$A$};
\node[vcol, minimum size=0.29cm] at (2.000,1.000) {$B$};
\node[vreg, minimum size=0.29cm] at (3.000,1.000) {$C$};
\node[vreg, dot] at (4.000,1.000) {};
\node[vreg, dot] at (0.000,2.000) {};
\node[vreg, dot] at (1.000,2.000) {};
\node[vreg, dot] at (2.000,2.000) {};
\node[vreg, dot] at (3.000,2.000) {};
\node[vreg, dot] at (4.000,2.000) {};

\end{tikzpicture}
}\caption{Rotação a partir de $c$.}\end{subfigure}\hfill
\begin{subfigure}[b]{0.32\linewidth}\centering\resizebox{\linewidth}{!}{\begin{tikzpicture}[scale=1.700, every node/.style={font=\scriptsize}]
\clip (0.450,0.300) rectangle (3.550,1.700);
\filldraw[cefill] (0.000,0.000) -- (1.000,0.000) -- (1.000,1.000) -- (0.000,1.000) -- cycle;
\filldraw[cefill] (1.000,0.000) -- (2.000,0.000) -- (2.000,0.800) -- (1.000,1.000) -- cycle;
\filldraw[cefill] (2.000,0.000) -- (3.000,0.000) -- (3.000,1.000) -- (2.000,0.800) -- cycle;
\filldraw[cefill] (3.000,0.000) -- (4.000,0.000) -- (4.000,1.000) -- (3.000,1.000) -- cycle;
\filldraw[cefill] (0.000,1.000) -- (1.000,1.000) -- (1.000,2.000) -- (0.000,2.000) -- cycle;
\filldraw[cefill] (1.000,1.000) -- (2.000,1.200) -- (2.000,2.000) -- (1.000,2.000) -- cycle;
\filldraw[cefill] (2.000,1.200) -- (3.000,1.000) -- (3.000,2.000) -- (2.000,2.000) -- cycle;
\filldraw[cefill] (3.000,1.000) -- (4.000,1.000) -- (4.000,2.000) -- (3.000,2.000) -- cycle;
\filldraw[iefill] (1.054,0.941) -- (2.000,0.800) -- (2.000,1.200) -- (1.054,1.059) -- cycle;
\filldraw[iefill] (2.000,0.800) -- (2.946,0.941) -- (2.946,1.059) -- (2.000,1.200) -- cycle;
\draw[red, line width=1.1pt] (1.054,1.059) -- (1.054,0.941);
\draw[red, line width=1.1pt] (2.946,0.941) -- (2.946,1.059);
\node[vreg, dot] at (0.000,0.000) {};
\node[vreg, dot] at (1.000,0.000) {};
\node[vreg, dot] at (2.000,0.000) {};
\node[vreg, dot] at (3.000,0.000) {};
\node[vreg, dot] at (4.000,0.000) {};
\node[vreg, dot] at (0.000,1.000) {};
\node[vcol, minimum size=0.29cm] at (1.000,1.000) {$A$};
\node[vie, minimum size=0.29cm] at (2.000,1.200) {$B$};
\node[vcol, minimum size=0.29cm] at (3.000,1.000) {$C$};
\node[vreg, dot] at (4.000,1.000) {};
\node[vreg, dot] at (0.000,2.000) {};
\node[vreg, dot] at (1.000,2.000) {};
\node[vreg, dot] at (2.000,2.000) {};
\node[vreg, dot] at (3.000,2.000) {};
\node[vreg, dot] at (4.000,2.000) {};
\node[vie, minimum size=0.29cm] at (2.000,0.800) {$B_{1}$};

\end{tikzpicture}
}\caption{Depois.}\end{subfigure}
\caption{Expand Edge em $B$: a aresta $BC$ continua a trinca $AB$, e o elemento de interface em $AB$ não toca os
elementos de $BC$. (b) Ordem em que \code{reindexElementsCCW} visita os cantos: 1 é a aresta colapsada $c$ (em $I$),
2 e 3 são os elementos abaixo da trinca; a rotação para ao encontrar $he$ (a half-edge de $BC$ que parte de $B$). (c)
Os cantos visitados passam a usar $B_1$; $A$ e $C$ continuam colapsados.}
\label{fig:expand}
\end{figure}

\subsection{Open Hole}
\label{sec:hole}

\textbf{Quando:} $v$ já pertence a um elemento de interface, mas não tem aresta colapsada. Isso acontece a partir da
terceira aresta em um vértice interno: depois do Split/Expand o vértice original já virou duas cópias, $v$ e $w$, e os
dois elementos de interface $J$ e $K$ que passam por ele estão colados pelos seus lados curtos (um vai de $v$ para $w$, o
outro de $w$ para $v$).

\textbf{O que faz} (\cref{alg:hole}): com três ou mais elementos de interface chegando ao mesmo ponto, os lados curtos não
conseguem mais fechar a vizinhança, e aparece um buraco. O operador cria $n$ e faz duas travessias a partir de $he$:
\begin{enumerate}
\item no sentido \emph{anti-horário}, troca $v$ por $n$ em todos os elementos até encontrar o primeiro elemento de
interface $J$; troca também o canto de $J$ e ``descola'' o lado curto de $J$ que chega a esse canto ($\Op\gets\BORDER$);
\item no sentido \emph{horário} (regra $h\leftarrow\nxt(\Op[h])$), apenas procura o primeiro elemento de interface $K$ e
descola o lado curto de $K$ que sai de $v$.
\end{enumerate}
O lado curto do novo elemento também fica na borda. Na primeira vez, os lados descolados de $J$ e $K$ eram justamente o
par que unia as duas cópias, e o resultado é um buraco triangular $(v,n,w)$; nas vezes seguintes o buraco já existe, os
lados encontrados já são de borda e o buraco ganha um vértice (\cref{fig:hole}).

\textbf{Vértice de borda.} O Open Hole também é o operador de um vértice de borda (externa ou de furo) que já foi aberto,
porque esse vértice está em um elemento de interface e não tem aresta colapsada. A estrela de um vértice de borda vai de
borda a borda, e por isso cada laço pode terminar em uma de duas ``paredes'': um elemento de interface, como num vértice
interno, ou a própria borda da malha. Quando um laço chega à borda, não há canto de elemento de interface a mover nem
lado curto a descolar: o buraco está aberto para o exterior, e aquele lado faz o mesmo que o Canônico de borda (o setor
entre a aresta nova e a borda passa para a cópia nova). Ao menos um dos lados sempre encontra um elemento de interface,
já que o vértice pertence a um; uma asserção verifica isso. Esse caso era o defeito I1 até a v0.4
(\cref{sec:falha-borda}).

\begin{algorithm}[ht]
\caption{\code{openHole}$(he, x)$}
\label{alg:hole}
\begin{algorithmic}[1]
\State $v\gets\V[he]$; \quad $n\gets$ \code{duplicatePoint}$(v)$; \quad $h\gets he$
\Repeat \Comment{laço 1: anti-horário, reindexando}
  \State $\V[h]\gets n$; \quad $h\gets\Op[\prv(h)]$
\Until{$h=\BORDER$ \textbf{ou} \code{\_isInterfaceElement}$[\mathrm{elem}(h)]$}
\State $\mathit{achouAH}\gets(h\ne\BORDER)$
\If{$\mathit{achouAH}$} $\V[h]\gets n$; \quad $\Op[\prv(h)]\gets\BORDER$ \Comment{canto e lado curto de $J$}
\EndIf
\State $h\gets he$
\Repeat \Comment{laço 2: horário, só procura}
  \State $o\gets\Op[h]$; \quad $h\gets$ (\textbf{se} $o=\BORDER$ \textbf{então} \BORDER{} \textbf{senão} $\nxt(o)$)
\Until{$h=\BORDER$ \textbf{ou} \code{\_isInterfaceElement}$[\mathrm{elem}(h)]$}
\State $\mathit{achouH}\gets(h\ne\BORDER)$
\If{$\mathit{achouH}$} $\Op[h]\gets\BORDER$ \Comment{lado curto de $K$}
\EndIf
\State \textbf{assert}$(\mathit{achouAH}$ \textbf{ou} $\mathit{achouH})$ \Comment{o vértice está em algum EI}
\State $\V[x]\gets n$; \quad $\V[\nxt(x)]\gets v$; \quad $\Op[x]\gets\BORDER$
\State marca $v$ e $n$ como vértices de EI
\end{algorithmic}
\end{algorithm}

\begin{figure}[ht]
\centering
\begin{subfigure}[b]{0.24\linewidth}\centering\resizebox{\linewidth}{!}{\begin{tikzpicture}[scale=1.900, every node/.style={font=\scriptsize}]
\clip (0.850,0.850) rectangle (3.150,3.150);
\filldraw[cefill] (0.000,0.000) -- (1.000,0.000) -- (1.000,1.000) -- (0.000,1.000) -- cycle;
\filldraw[cefill] (1.000,0.000) -- (2.000,0.000) -- (2.000,1.000) -- (1.000,1.000) -- cycle;
\filldraw[cefill] (2.000,0.000) -- (3.000,0.000) -- (3.000,1.000) -- (2.000,1.000) -- cycle;
\filldraw[cefill] (3.000,0.000) -- (4.000,0.000) -- (4.000,1.000) -- (3.000,1.000) -- cycle;
\filldraw[cefill] (0.000,1.000) -- (1.000,1.000) -- (1.000,2.000) -- (0.000,2.000) -- cycle;
\filldraw[cefill] (1.000,1.000) -- (2.000,1.000) -- (2.000,2.000) -- (1.000,2.000) -- cycle;
\filldraw[cefill] (2.000,1.000) -- (3.000,1.000) -- (3.000,2.000) -- (2.000,2.000) -- cycle;
\filldraw[cefill] (3.000,1.000) -- (4.000,1.000) -- (4.000,2.000) -- (3.000,2.000) -- cycle;
\filldraw[cefill] (0.000,2.000) -- (1.000,2.000) -- (1.000,3.000) -- (0.000,3.000) -- cycle;
\filldraw[cefill] (1.000,2.000) -- (2.000,2.000) -- (2.000,3.000) -- (1.000,3.000) -- cycle;
\filldraw[cefill] (2.000,2.000) -- (3.000,2.000) -- (3.000,3.000) -- (2.000,3.000) -- cycle;
\filldraw[cefill] (3.000,2.000) -- (4.000,2.000) -- (4.000,3.000) -- (3.000,3.000) -- cycle;
\filldraw[cefill] (0.000,3.000) -- (1.000,3.000) -- (1.000,4.000) -- (0.000,4.000) -- cycle;
\filldraw[cefill] (1.000,3.000) -- (2.000,3.000) -- (2.000,4.000) -- (1.000,4.000) -- cycle;
\filldraw[cefill] (2.000,3.000) -- (3.000,3.000) -- (3.000,4.000) -- (2.000,4.000) -- cycle;
\filldraw[cefill] (3.000,3.000) -- (4.000,3.000) -- (4.000,4.000) -- (3.000,4.000) -- cycle;
\filldraw[iefill] (1.074,1.926) -- (1.926,1.926) -- (1.926,2.074) -- (1.074,2.074) -- cycle;
\draw[red, line width=1.1pt] (1.926,1.926) -- (1.926,2.074);
\draw[red, line width=1.1pt] (1.074,2.074) -- (1.074,1.926);
\draw[curedge] (2.000,2.000) -- (3.000,2.000);
\node[edgenum] at (2.500,2.160) {2};
\draw[seledge] (2.000,2.000) -- (2.000,3.000);
\node[edgenum] at (1.840,2.500) {3};
\draw[seledge] (2.000,1.000) -- (2.000,2.000);
\node[edgenum] at (1.840,1.500) {4};
\node[vreg, font=\tiny\bfseries, dot] at (0.000,0.000) {};
\node[vreg, font=\tiny\bfseries, dot] at (1.000,0.000) {};
\node[vreg, font=\tiny\bfseries, dot] at (2.000,0.000) {};
\node[vreg, font=\tiny\bfseries, dot] at (3.000,0.000) {};
\node[vreg, font=\tiny\bfseries, dot] at (4.000,0.000) {};
\node[vreg, font=\tiny\bfseries, dot] at (0.000,1.000) {};
\node[vreg, font=\tiny\bfseries, dot] at (1.000,1.000) {};
\node[vreg, font=\tiny\bfseries, dot] at (2.000,1.000) {};
\node[vreg, font=\tiny\bfseries, dot] at (3.000,1.000) {};
\node[vreg, font=\tiny\bfseries, dot] at (4.000,1.000) {};
\node[vreg, font=\tiny\bfseries, dot] at (0.000,2.000) {};
\node[vcol, font=\tiny\bfseries, dot] at (1.000,2.000) {};
\node[vcol, font=\tiny\bfseries, minimum size=0.24cm] at (2.000,2.000) {$V$};
\node[vreg, font=\tiny\bfseries, dot] at (3.000,2.000) {};
\node[vreg, font=\tiny\bfseries, dot] at (4.000,2.000) {};
\node[vreg, font=\tiny\bfseries, dot] at (0.000,3.000) {};
\node[vreg, font=\tiny\bfseries, dot] at (1.000,3.000) {};
\node[vreg, font=\tiny\bfseries, dot] at (2.000,3.000) {};
\node[vreg, font=\tiny\bfseries, dot] at (3.000,3.000) {};
\node[vreg, font=\tiny\bfseries, dot] at (4.000,3.000) {};
\node[vreg, font=\tiny\bfseries, dot] at (0.000,4.000) {};
\node[vreg, font=\tiny\bfseries, dot] at (1.000,4.000) {};
\node[vreg, font=\tiny\bfseries, dot] at (2.000,4.000) {};
\node[vreg, font=\tiny\bfseries, dot] at (3.000,4.000) {};
\node[vreg, font=\tiny\bfseries, dot] at (4.000,4.000) {};

\end{tikzpicture}
}\caption{Aresta 1: Canônico.}\end{subfigure}\hfill
\begin{subfigure}[b]{0.24\linewidth}\centering\resizebox{\linewidth}{!}{\begin{tikzpicture}[scale=1.900, every node/.style={font=\scriptsize}]
\clip (0.850,0.850) rectangle (3.150,3.150);
\filldraw[cefill] (0.000,0.000) -- (1.000,0.000) -- (1.000,1.000) -- (0.000,1.000) -- cycle;
\filldraw[cefill] (1.000,0.000) -- (2.000,0.000) -- (2.000,1.000) -- (1.000,1.000) -- cycle;
\filldraw[cefill] (2.000,0.000) -- (3.000,0.000) -- (3.000,1.000) -- (2.000,1.000) -- cycle;
\filldraw[cefill] (3.000,0.000) -- (4.000,0.000) -- (4.000,1.000) -- (3.000,1.000) -- cycle;
\filldraw[cefill] (0.000,1.000) -- (1.000,1.000) -- (1.000,2.000) -- (0.000,2.000) -- cycle;
\filldraw[cefill] (1.000,1.000) -- (2.000,1.000) -- (2.000,1.740) -- (1.000,2.000) -- cycle;
\filldraw[cefill] (2.000,1.000) -- (3.000,1.000) -- (3.000,2.000) -- (2.000,1.740) -- cycle;
\filldraw[cefill] (3.000,1.000) -- (4.000,1.000) -- (4.000,2.000) -- (3.000,2.000) -- cycle;
\filldraw[cefill] (0.000,2.000) -- (1.000,2.000) -- (1.000,3.000) -- (0.000,3.000) -- cycle;
\filldraw[cefill] (1.000,2.000) -- (2.000,2.260) -- (2.000,3.000) -- (1.000,3.000) -- cycle;
\filldraw[cefill] (2.000,2.260) -- (3.000,2.000) -- (3.000,3.000) -- (2.000,3.000) -- cycle;
\filldraw[cefill] (3.000,2.000) -- (4.000,2.000) -- (4.000,3.000) -- (3.000,3.000) -- cycle;
\filldraw[cefill] (0.000,3.000) -- (1.000,3.000) -- (1.000,4.000) -- (0.000,4.000) -- cycle;
\filldraw[cefill] (1.000,3.000) -- (2.000,3.000) -- (2.000,4.000) -- (1.000,4.000) -- cycle;
\filldraw[cefill] (2.000,3.000) -- (3.000,3.000) -- (3.000,4.000) -- (2.000,4.000) -- cycle;
\filldraw[cefill] (3.000,3.000) -- (4.000,3.000) -- (4.000,4.000) -- (3.000,4.000) -- cycle;
\filldraw[iefill] (1.069,1.922) -- (2.000,1.740) -- (2.000,2.260) -- (1.069,2.078) -- cycle;
\filldraw[iefill] (2.000,1.740) -- (2.931,1.922) -- (2.931,2.078) -- (2.000,2.260) -- cycle;
\draw[red, line width=1.1pt] (1.069,2.078) -- (1.069,1.922);
\draw[red, line width=1.1pt] (2.931,1.922) -- (2.931,2.078);
\draw[curedge] (2.000,2.260) -- (2.000,3.000);
\node[edgenum] at (1.840,2.630) {3};
\draw[seledge] (2.000,1.000) -- (2.000,1.740);
\node[edgenum] at (1.840,1.370) {4};
\node[vreg, font=\tiny\bfseries, dot] at (0.000,0.000) {};
\node[vreg, font=\tiny\bfseries, dot] at (1.000,0.000) {};
\node[vreg, font=\tiny\bfseries, dot] at (2.000,0.000) {};
\node[vreg, font=\tiny\bfseries, dot] at (3.000,0.000) {};
\node[vreg, font=\tiny\bfseries, dot] at (4.000,0.000) {};
\node[vreg, font=\tiny\bfseries, dot] at (0.000,1.000) {};
\node[vreg, font=\tiny\bfseries, dot] at (1.000,1.000) {};
\node[vreg, font=\tiny\bfseries, dot] at (2.000,1.000) {};
\node[vreg, font=\tiny\bfseries, dot] at (3.000,1.000) {};
\node[vreg, font=\tiny\bfseries, dot] at (4.000,1.000) {};
\node[vreg, font=\tiny\bfseries, dot] at (0.000,2.000) {};
\node[vcol, font=\tiny\bfseries, dot] at (1.000,2.000) {};
\node[vie, font=\tiny\bfseries, minimum size=0.24cm] at (2.000,2.260) {$V$};
\node[vcol, font=\tiny\bfseries, dot] at (3.000,2.000) {};
\node[vreg, font=\tiny\bfseries, dot] at (4.000,2.000) {};
\node[vreg, font=\tiny\bfseries, dot] at (0.000,3.000) {};
\node[vreg, font=\tiny\bfseries, dot] at (1.000,3.000) {};
\node[vreg, font=\tiny\bfseries, dot] at (2.000,3.000) {};
\node[vreg, font=\tiny\bfseries, dot] at (3.000,3.000) {};
\node[vreg, font=\tiny\bfseries, dot] at (4.000,3.000) {};
\node[vreg, font=\tiny\bfseries, dot] at (0.000,4.000) {};
\node[vreg, font=\tiny\bfseries, dot] at (1.000,4.000) {};
\node[vreg, font=\tiny\bfseries, dot] at (2.000,4.000) {};
\node[vreg, font=\tiny\bfseries, dot] at (3.000,4.000) {};
\node[vreg, font=\tiny\bfseries, dot] at (4.000,4.000) {};
\node[vie, font=\tiny\bfseries, minimum size=0.24cm] at (2.000,1.740) {$V_{1}$};

\end{tikzpicture}
}\caption{Aresta 2: Expand.}\end{subfigure}\hfill
\begin{subfigure}[b]{0.24\linewidth}\centering\resizebox{\linewidth}{!}{\begin{tikzpicture}[scale=1.900, every node/.style={font=\scriptsize}]
\clip (0.850,0.850) rectangle (3.150,3.150);
\fill[holefill] (2.000,1.740) -- (1.816,2.184) -- (2.184,2.184) -- cycle;
\filldraw[cefill] (0.000,0.000) -- (1.000,0.000) -- (1.000,1.000) -- (0.000,1.000) -- cycle;
\filldraw[cefill] (1.000,0.000) -- (2.000,0.000) -- (2.000,1.000) -- (1.000,1.000) -- cycle;
\filldraw[cefill] (2.000,0.000) -- (3.000,0.000) -- (3.000,1.000) -- (2.000,1.000) -- cycle;
\filldraw[cefill] (3.000,0.000) -- (4.000,0.000) -- (4.000,1.000) -- (3.000,1.000) -- cycle;
\filldraw[cefill] (0.000,1.000) -- (1.000,1.000) -- (1.000,2.000) -- (0.000,2.000) -- cycle;
\filldraw[cefill] (1.000,1.000) -- (2.000,1.000) -- (2.000,1.740) -- (1.000,2.000) -- cycle;
\filldraw[cefill] (2.000,1.000) -- (3.000,1.000) -- (3.000,2.000) -- (2.000,1.740) -- cycle;
\filldraw[cefill] (3.000,1.000) -- (4.000,1.000) -- (4.000,2.000) -- (3.000,2.000) -- cycle;
\filldraw[cefill] (0.000,2.000) -- (1.000,2.000) -- (1.000,3.000) -- (0.000,3.000) -- cycle;
\filldraw[cefill] (1.000,2.000) -- (1.816,2.184) -- (2.000,3.000) -- (1.000,3.000) -- cycle;
\filldraw[cefill] (2.184,2.184) -- (3.000,2.000) -- (3.000,3.000) -- (2.000,3.000) -- cycle;
\filldraw[cefill] (3.000,2.000) -- (4.000,2.000) -- (4.000,3.000) -- (3.000,3.000) -- cycle;
\filldraw[cefill] (0.000,3.000) -- (1.000,3.000) -- (1.000,4.000) -- (0.000,4.000) -- cycle;
\filldraw[cefill] (1.000,3.000) -- (2.000,3.000) -- (2.000,4.000) -- (1.000,4.000) -- cycle;
\filldraw[cefill] (2.000,3.000) -- (3.000,3.000) -- (3.000,4.000) -- (2.000,4.000) -- cycle;
\filldraw[cefill] (3.000,3.000) -- (4.000,3.000) -- (4.000,4.000) -- (3.000,4.000) -- cycle;
\filldraw[iefill] (1.069,1.922) -- (2.000,1.740) -- (1.816,2.184) -- (1.066,2.080) -- cycle;
\filldraw[iefill] (2.000,1.740) -- (2.931,1.922) -- (2.934,2.080) -- (2.184,2.184) -- cycle;
\filldraw[iefill] (2.184,2.184) -- (2.080,2.934) -- (1.920,2.934) -- (1.816,2.184) -- cycle;
\draw[red, line width=1.1pt] (1.066,2.080) -- (1.069,1.922);
\draw[red, line width=1.1pt] (2.931,1.922) -- (2.934,2.080);
\draw[red, line width=1.1pt] (2.080,2.934) -- (1.920,2.934);
\draw[curedge] (2.000,1.000) -- (2.000,1.740);
\node[edgenum] at (1.840,1.370) {4};
\node[vreg, font=\tiny\bfseries, dot] at (0.000,0.000) {};
\node[vreg, font=\tiny\bfseries, dot] at (1.000,0.000) {};
\node[vreg, font=\tiny\bfseries, dot] at (2.000,0.000) {};
\node[vreg, font=\tiny\bfseries, dot] at (3.000,0.000) {};
\node[vreg, font=\tiny\bfseries, dot] at (4.000,0.000) {};
\node[vreg, font=\tiny\bfseries, dot] at (0.000,1.000) {};
\node[vreg, font=\tiny\bfseries, dot] at (1.000,1.000) {};
\node[vreg, font=\tiny\bfseries, dot] at (2.000,1.000) {};
\node[vreg, font=\tiny\bfseries, dot] at (3.000,1.000) {};
\node[vreg, font=\tiny\bfseries, dot] at (4.000,1.000) {};
\node[vreg, font=\tiny\bfseries, dot] at (0.000,2.000) {};
\node[vcol, font=\tiny\bfseries, dot] at (1.000,2.000) {};
\node[vie, font=\tiny\bfseries, minimum size=0.24cm] at (2.184,2.184) {$V$};
\node[vcol, font=\tiny\bfseries, dot] at (3.000,2.000) {};
\node[vreg, font=\tiny\bfseries, dot] at (4.000,2.000) {};
\node[vreg, font=\tiny\bfseries, dot] at (0.000,3.000) {};
\node[vreg, font=\tiny\bfseries, dot] at (1.000,3.000) {};
\node[vcol, font=\tiny\bfseries, dot] at (2.000,3.000) {};
\node[vreg, font=\tiny\bfseries, dot] at (3.000,3.000) {};
\node[vreg, font=\tiny\bfseries, dot] at (4.000,3.000) {};
\node[vreg, font=\tiny\bfseries, dot] at (0.000,4.000) {};
\node[vreg, font=\tiny\bfseries, dot] at (1.000,4.000) {};
\node[vreg, font=\tiny\bfseries, dot] at (2.000,4.000) {};
\node[vreg, font=\tiny\bfseries, dot] at (3.000,4.000) {};
\node[vreg, font=\tiny\bfseries, dot] at (4.000,4.000) {};
\node[vie, font=\tiny\bfseries, minimum size=0.24cm] at (2.000,1.740) {$V_{1}$};
\node[vie, font=\tiny\bfseries, minimum size=0.24cm] at (1.816,2.184) {$V_{2}$};

\end{tikzpicture}
}\caption{Aresta 3: Open Hole.}\end{subfigure}\hfill
\begin{subfigure}[b]{0.24\linewidth}\centering\resizebox{\linewidth}{!}{\begin{tikzpicture}[scale=1.900, every node/.style={font=\scriptsize}]
\clip (0.850,0.850) rectangle (3.150,3.150);
\fill[holefill] (1.816,1.816) -- (1.816,2.184) -- (2.184,2.184) -- (2.184,1.816) -- cycle;
\filldraw[cefill] (0.000,0.000) -- (1.000,0.000) -- (1.000,1.000) -- (0.000,1.000) -- cycle;
\filldraw[cefill] (1.000,0.000) -- (2.000,0.000) -- (2.000,1.000) -- (1.000,1.000) -- cycle;
\filldraw[cefill] (2.000,0.000) -- (3.000,0.000) -- (3.000,1.000) -- (2.000,1.000) -- cycle;
\filldraw[cefill] (3.000,0.000) -- (4.000,0.000) -- (4.000,1.000) -- (3.000,1.000) -- cycle;
\filldraw[cefill] (0.000,1.000) -- (1.000,1.000) -- (1.000,2.000) -- (0.000,2.000) -- cycle;
\filldraw[cefill] (1.000,1.000) -- (2.000,1.000) -- (1.816,1.816) -- (1.000,2.000) -- cycle;
\filldraw[cefill] (2.000,1.000) -- (3.000,1.000) -- (3.000,2.000) -- (2.184,1.816) -- cycle;
\filldraw[cefill] (3.000,1.000) -- (4.000,1.000) -- (4.000,2.000) -- (3.000,2.000) -- cycle;
\filldraw[cefill] (0.000,2.000) -- (1.000,2.000) -- (1.000,3.000) -- (0.000,3.000) -- cycle;
\filldraw[cefill] (1.000,2.000) -- (1.816,2.184) -- (2.000,3.000) -- (1.000,3.000) -- cycle;
\filldraw[cefill] (2.184,2.184) -- (3.000,2.000) -- (3.000,3.000) -- (2.000,3.000) -- cycle;
\filldraw[cefill] (3.000,2.000) -- (4.000,2.000) -- (4.000,3.000) -- (3.000,3.000) -- cycle;
\filldraw[cefill] (0.000,3.000) -- (1.000,3.000) -- (1.000,4.000) -- (0.000,4.000) -- cycle;
\filldraw[cefill] (1.000,3.000) -- (2.000,3.000) -- (2.000,4.000) -- (1.000,4.000) -- cycle;
\filldraw[cefill] (2.000,3.000) -- (3.000,3.000) -- (3.000,4.000) -- (2.000,4.000) -- cycle;
\filldraw[cefill] (3.000,3.000) -- (4.000,3.000) -- (4.000,4.000) -- (3.000,4.000) -- cycle;
\filldraw[iefill] (1.066,1.920) -- (1.816,1.816) -- (1.816,2.184) -- (1.066,2.080) -- cycle;
\filldraw[iefill] (2.184,1.816) -- (2.934,1.920) -- (2.934,2.080) -- (2.184,2.184) -- cycle;
\filldraw[iefill] (2.184,2.184) -- (2.080,2.934) -- (1.920,2.934) -- (1.816,2.184) -- cycle;
\filldraw[iefill] (2.080,1.066) -- (2.184,1.816) -- (1.816,1.816) -- (1.920,1.066) -- cycle;
\draw[red, line width=1.1pt] (1.066,2.080) -- (1.066,1.920);
\draw[red, line width=1.1pt] (2.934,1.920) -- (2.934,2.080);
\draw[red, line width=1.1pt] (2.080,2.934) -- (1.920,2.934);
\draw[red, line width=1.1pt] (1.920,1.066) -- (2.080,1.066);
\node[vreg, font=\tiny\bfseries, dot] at (0.000,0.000) {};
\node[vreg, font=\tiny\bfseries, dot] at (1.000,0.000) {};
\node[vreg, font=\tiny\bfseries, dot] at (2.000,0.000) {};
\node[vreg, font=\tiny\bfseries, dot] at (3.000,0.000) {};
\node[vreg, font=\tiny\bfseries, dot] at (4.000,0.000) {};
\node[vreg, font=\tiny\bfseries, dot] at (0.000,1.000) {};
\node[vreg, font=\tiny\bfseries, dot] at (1.000,1.000) {};
\node[vcol, font=\tiny\bfseries, dot] at (2.000,1.000) {};
\node[vreg, font=\tiny\bfseries, dot] at (3.000,1.000) {};
\node[vreg, font=\tiny\bfseries, dot] at (4.000,1.000) {};
\node[vreg, font=\tiny\bfseries, dot] at (0.000,2.000) {};
\node[vcol, font=\tiny\bfseries, dot] at (1.000,2.000) {};
\node[vie, font=\tiny\bfseries, minimum size=0.24cm] at (2.184,2.184) {$V$};
\node[vcol, font=\tiny\bfseries, dot] at (3.000,2.000) {};
\node[vreg, font=\tiny\bfseries, dot] at (4.000,2.000) {};
\node[vreg, font=\tiny\bfseries, dot] at (0.000,3.000) {};
\node[vreg, font=\tiny\bfseries, dot] at (1.000,3.000) {};
\node[vcol, font=\tiny\bfseries, dot] at (2.000,3.000) {};
\node[vreg, font=\tiny\bfseries, dot] at (3.000,3.000) {};
\node[vreg, font=\tiny\bfseries, dot] at (4.000,3.000) {};
\node[vreg, font=\tiny\bfseries, dot] at (0.000,4.000) {};
\node[vreg, font=\tiny\bfseries, dot] at (1.000,4.000) {};
\node[vreg, font=\tiny\bfseries, dot] at (2.000,4.000) {};
\node[vreg, font=\tiny\bfseries, dot] at (3.000,4.000) {};
\node[vreg, font=\tiny\bfseries, dot] at (4.000,4.000) {};
\node[vie, font=\tiny\bfseries, minimum size=0.24cm] at (1.816,1.816) {$V_{1}$};
\node[vie, font=\tiny\bfseries, minimum size=0.24cm] at (1.816,2.184) {$V_{2}$};
\node[vie, font=\tiny\bfseries, minimum size=0.24cm] at (2.184,1.816) {$V_{3}$};

\end{tikzpicture}
}\caption{Aresta 4: Open Hole.}\end{subfigure}
\caption{Quatro arestas chegando ao vértice interno $V$ (detalhe do centro de uma malha $4\times4$; a próxima aresta a
inserir está em vermelho grosso). (a) A primeira deixa $V$ colapsado. (b) A segunda, alinhada, divide $V$ em duas cópias
coladas pelos lados curtos. (c) A terceira descola esses lados: buraco triangular (hachurado). (d) A quarta aumenta o
buraco para um quadrilátero. Com $k$ arestas, o buraco tem $k$ vértices.}
\label{fig:hole}
\end{figure}

\subsection{Relações entre os operadores}
\label{sec:relacoes}

Olhando os quatro operadores lado a lado, todos fazem a mesma coisa em essência: \emph{dividem um leque da estrela de $v$
entre duas ``paredes'', dão um vértice novo a um dos lados e decidem o que fazer com os lados curtos}. O que muda é a
parede onde a rotação começa e onde ela para:

\begin{center}\small
\begin{tabular}{llll}
\toprule
Operador & começa em & para em & lados curtos\\
\midrule
Canônico (borda) & aresta nova & borda da malha & novo lado na borda\\
Expand Edge & aresta colapsada $c$ & aresta nova & $c$ é colada ao novo lado\\
Split Element & aresta colapsada $c$ & aresta nova (após 1 elemento) & $c$ é colada ao novo lado\\
Open Hole & aresta nova & primeiro elemento de interface (ou a borda) & descolados (buraco)\\
\bottomrule
\end{tabular}
\end{center}

Uma consequência interessante: o \textbf{Split Element é um caso particular do Expand Edge}. Na hipótese $x_3$, a rotação
do Expand a partir de $c$ visita exatamente o canto de $c$ e o canto $\nxt(b)$ de $S$ e para em $he$ --- as mesmas
alterações do Split. Na hipótese $x_1$, a rotação do Expand vai pelo outro lado e reindexa todos os elementos da estrela
\emph{exceto} $S$; o resultado é a mesma divisão da estrela, só com os nomes $v$ e $n$ trocados. Verificamos isso
experimentalmente: substituindo o Split pelo Expand em 791 ocorrências (400 malhas aleatórias de triângulos,
quadriláteros, híbridas e com polígonos), a topologia final foi idêntica em todos os casos (\cref{sec:verificacao}).
Desde a v0.5 essa equivalência é um teste automatizado, em malhas aleatórias dos cinco tipos (\cref{sec:testes}). O
Split Element é, portanto, uma \emph{otimização}: custo $O(\log n_c)$ em vez de $O(d\log n_c)$, movendo sempre o lado
menor. Isso não é um defeito; mas é uma informação útil para a extensão 3D,
onde talvez valha começar pela versão geral (rotação) e só depois especializar.

\subsection{Contagem prévia de vértices}
\label{sec:contagem}

Para reservar memória de uma vez, \code{computeNumberOfNewVertices} calcula, antes de alterar a malha, quantos vértices
a chamada vai criar. A contagem tem duas etapas:
\begin{enumerate}
\item \textbf{Seleção.} Percorre as arestas do lote e descarta as mesmas que \code{insertInterfaceElements} vai
rejeitar, com os mesmos testes e na mesma ordem: índice que não é half-edge da malha ($h\ge Of[n_c]$), aresta de borda,
aresta de elemento de interface e aresta repetida no lote (inclusive quando vem uma vez por cada uma de suas
half-edges; a aresta é identificada pela menor das duas). Para cada vértice tocado pelas arestas que sobram, guarda o
número $n_v$ dessas arestas e uma half-edge que parte dele.
\item \textbf{Regra por vértice.} Aplica, uma vez por vértice,
\[
\text{novos}(v)=
\begin{cases}
n_v-1, & v \text{ fora de EI e interno (o primeiro contato só colapsa)},\\
n_v, & v \text{ fora de EI e na borda (o primeiro contato já duplica)},\\
n_v, & v \text{ já em algum EI (todo contato duplica)}.
\end{cases}
\]
\end{enumerate}
Como a regra é aplicada por vértice, e não por aresta, o teste de borda (\code{isBorder}, que gira em torno do vértice)
roda uma vez por vértice distinto, e um vértice tocado por várias arestas do lote é contado corretamente. Por descartar
as arestas de elementos de interface antes do teste de borda, a contagem também deixou de chegar ao caminho em que
\code{isBorder} não termina (I2). O resultado é exatamente o número de vértices criados, em todos os testes
(\cref{sec:verificacao}); até a v0.3 a contagem era um limite superior (defeito I5, corrigido). Num vértice de borda com
duas ou mais arestas, ela dá o valor correto ($n_v$), uma cópia por aresta (Open Hole na borda). Os elementos
são reservados à razão de um por aresta do lote.
